% --- Archon physics formalization source begin ---
% archon:physics
% archon:covers PhyXMiniProblems/problem_phyx_mini_0504.lean
% archon:source-report reports/phyx_mini/problem_phyx_mini_0504.source.json

\chapter{Physics problem phyx\_mini\_0504}
\label{ch:PhyXMiniProblems_problem_phyx_mini_0504}

\paragraph{Problem source.}
\#\# Physical scenario

Figure shows a pulse train. The period of the pulse train is \textbackslash{}( T = 2\textbackslash{}Delta t \textbackslash{}), where \textbackslash{}( \textbackslash{}Delta t \textbackslash{}) is the duration of each pulse.

\#\# Question

What is the maximum pulse-transmission rate (pulses per second) through an electronics system with a \textbackslash{}( 200 \textbackslash{}text\{ kHz\} \textbackslash{}) bandwidth? (This is the bandwidth allotted to each FM radio station.)

\#\# Answer choices

- A: A: \textbackslash{}( 1.0 \textbackslash{}times 10\textasciicircum{}3 \textbackslash{}text\{ pulses/s\} \textbackslash{})

- B: B: \textbackslash{}( 1.0 \textbackslash{}times 10\textasciicircum{}7 \textbackslash{}text\{ pulses/s\} \textbackslash{})

- C: C: \textbackslash{}( 1.0 \textbackslash{}times 10\textasciicircum{}5 \textbackslash{}text\{ pulses/s\} \textbackslash{})

- D: D: \textbackslash{}( 1.0 \textbackslash{}times 10\textasciicircum{}6 \textbackslash{}text\{ pulses/s\} \textbackslash{})

\#\# Figure caption (auxiliary; use the image as primary evidence)

The image depicts a periodic square wave signal. Here's a detailed breakdown:

- **Signal Shape**: The wave consists of a series of rectangular pulses, indicating a square wave.
- **Axes**: 
  - The horizontal axis is labeled "t" (representing time).
- **Pulse Details**: 
  - Each pulse has a width of \textbackslash{}(\textbackslash{}Delta t\textbackslash{}).
  - The period between the start of one pulse and the start of the next is labeled as \textbackslash{}(T = 2\textbackslash{}Delta t\textbackslash{}).
- **Labels**:
  - Above the waveform, the width of the pulse is marked with a double-headed arrow and labeled \textbackslash{}(\textbackslash{}Delta t\textbackslash{}).
  - Below the waveform, a section corresponding to one complete period is marked with a double-headed arrow and labeled "Period \textbackslash{}(T = 2\textbackslash{}Delta t\textbackslash{})".

The waveform is depicted in purple, while the horizontal axis is in black.

\#\# Dataset metadata

Quantum Phenomena; Numerical Reasoning

\paragraph{Recorded answer/context.}
C: C: \textbackslash{}( 1.0 \textbackslash{}times 10\textasciicircum{}5 \textbackslash{}text\{ pulses/s\} \textbackslash{})

\paragraph{Figure/image path.}
/root/proposal\_for\_physic/science-mango/phyx\_mini\_run/phyx\_data/test\_image/504.png

\paragraph{Formalization target.}
create a compiling Lean file with sorry bodies at `PhyXMiniProblems/problem\_phyx\_mini\_0504.lean`. The Lean declarations must preserve the physical quantities, dimensions or dimensional roles, figure labels, governing-law hypotheses, and final relation expressed by this problem.
Use Mathlib/Physlib names found through LeanExplore where available. If a physics API is missing, introduce faithful local abstractions rather than scalar placeholder aliases.

\begin{theorem}[Physics formalization target]
\label{thm:physics:phyx_mini_0504:target}
\lean{PhyXMiniProblems.ProblemPhyXMini0504.problem_phyx_mini_0504}
\uses{def:physics:phyx-mini-0504:phyxminiproblems-problemphyxmini0504-pulserateinpulsespersecond, def:physics:phyx-mini-0504:phyxminiproblems-problemphyxmini0504-bandwidthlimitedpulsetrain, def:physics:phyx-mini-0504:phyxminiproblems-problemphyxmini0504-matchesproblemstatementandfigure, def:physics:phyx-mini-0504:phyxminiproblems-problemphyxmini0504-hasphysicalpulsetrainparameters, def:physics:phyx-mini-0504:phyxminiproblems-problemphyxmini0504-satisfiesidealpulsebandwidthlaws, def:physics:phyx-mini-0504:phyxminiproblems-problemphyxmini0504-recordeddatasetanswer, def:physics:phyx-mini-0504:phyxminiproblems-problemphyxmini0504-matchesanswerchoice}
The assigned autoformalize agent should translate this physics problem into Lean declarations in the covered file, with theorem and lemma proof bodies written as `by sorry`.
\end{theorem}
\begin{proof}
This is an autoformalization task, not a proof task. Produce statements that can later be proved by the physics prover without weakening the physical model.
\end{proof}

% --- Archon declaration topology begin ---
\section{Declaration topology}
\begin{definition}[Duration Quantity]
  \label{def:physics:phyx-mini-0504:phyxminiproblems-problemphyxmini0504-durationquantity}
  \lean{PhyXMiniProblems.ProblemPhyXMini0504.DurationQuantity}
  A nonnegative physical duration, carrying the dimension of time
\end{definition}
\begin{proof}
  This is the declared model definition of Duration Quantity.
\end{proof}
\begin{definition}[Frequency Quantity]
  \label{def:physics:phyx-mini-0504:phyxminiproblems-problemphyxmini0504-frequencyquantity}
  \lean{PhyXMiniProblems.ProblemPhyXMini0504.FrequencyQuantity}
  A nonnegative physical frequency, carrying inverse-time dimension
\end{definition}
\begin{proof}
  This is the declared model definition of Frequency Quantity.
\end{proof}
\begin{definition}[duration Readout]
  \label{def:physics:phyx-mini-0504:phyxminiproblems-problemphyxmini0504-durationreadout}
  \lean{PhyXMiniProblems.ProblemPhyXMini0504.durationReadout}
  \uses{def:physics:phyx-mini-0504:phyxminiproblems-problemphyxmini0504-durationquantity}
  Read a physical duration in the selected time unit
\end{definition}
\begin{proof}
  This is the declared model definition of duration Readout.
\end{proof}
\begin{definition}[frequency Readout]
  \label{def:physics:phyx-mini-0504:phyxminiproblems-problemphyxmini0504-frequencyreadout}
  \lean{PhyXMiniProblems.ProblemPhyXMini0504.frequencyReadout}
  \uses{def:physics:phyx-mini-0504:phyxminiproblems-problemphyxmini0504-frequencyquantity}
  Read a physical inverse-time quantity in the inverse selected time unit
\end{definition}
\begin{proof}
  This is the declared model definition of frequency Readout.
\end{proof}
\begin{definition}[duration In Seconds]
  \label{def:physics:phyx-mini-0504:phyxminiproblems-problemphyxmini0504-durationinseconds}
  \lean{PhyXMiniProblems.ProblemPhyXMini0504.durationInSeconds}
  \uses{def:physics:phyx-mini-0504:phyxminiproblems-problemphyxmini0504-durationquantity, def:physics:phyx-mini-0504:phyxminiproblems-problemphyxmini0504-durationreadout}
  Duration readout in seconds
\end{definition}
\begin{proof}
  This is the declared model definition of duration In Seconds.
\end{proof}
\begin{definition}[frequency In Hertz]
  \label{def:physics:phyx-mini-0504:phyxminiproblems-problemphyxmini0504-frequencyinhertz}
  \lean{PhyXMiniProblems.ProblemPhyXMini0504.frequencyInHertz}
  \uses{def:physics:phyx-mini-0504:phyxminiproblems-problemphyxmini0504-frequencyquantity, def:physics:phyx-mini-0504:phyxminiproblems-problemphyxmini0504-frequencyreadout}
  Frequency readout in hertz, i.e. inverse seconds
\end{definition}
\begin{proof}
  This is the declared model definition of frequency In Hertz.
\end{proof}
\begin{definition}[bandwidth In Kilohertz]
  \label{def:physics:phyx-mini-0504:phyxminiproblems-problemphyxmini0504-bandwidthinkilohertz}
  \lean{PhyXMiniProblems.ProblemPhyXMini0504.bandwidthInKilohertz}
  \uses{def:physics:phyx-mini-0504:phyxminiproblems-problemphyxmini0504-frequencyquantity, def:physics:phyx-mini-0504:phyxminiproblems-problemphyxmini0504-frequencyreadout}
  Bandwidth readout in kilohertz, i.e. inverse milliseconds
\end{definition}
\begin{proof}
  This is the declared model definition of bandwidth In Kilohertz.
\end{proof}
\begin{definition}[pulse Rate In Pulses Per Second]
  \label{def:physics:phyx-mini-0504:phyxminiproblems-problemphyxmini0504-pulserateinpulsespersecond}
  \lean{PhyXMiniProblems.ProblemPhyXMini0504.pulseRateInPulsesPerSecond}
  \uses{def:physics:phyx-mini-0504:phyxminiproblems-problemphyxmini0504-frequencyquantity, def:physics:phyx-mini-0504:phyxminiproblems-problemphyxmini0504-frequencyinhertz}
  A pulse-rate readout in pulses per second; pulse count is dimensionless
\end{definition}
\begin{proof}
  This is the declared model definition of pulse Rate In Pulses Per Second.
\end{proof}
\begin{definition}[Pulse Train Shape]
  \label{def:physics:phyx-mini-0504:phyxminiproblems-problemphyxmini0504-pulsetrainshape}
  \lean{PhyXMiniProblems.ProblemPhyXMini0504.PulseTrainShape}
  The waveform shape visible in the primary image
\end{definition}
\begin{proof}
  This is the declared model definition of Pulse Train Shape.
\end{proof}
\begin{definition}[Figure Label Role]
  \label{def:physics:phyx-mini-0504:phyxminiproblems-problemphyxmini0504-figurelabelrole}
  \lean{PhyXMiniProblems.ProblemPhyXMini0504.FigureLabelRole}
  Roles of the labels printed in the primary image
\end{definition}
\begin{proof}
  This is the declared model definition of Figure Label Role.
\end{proof}
\begin{definition}[Figure Label]
  \label{def:physics:phyx-mini-0504:phyxminiproblems-problemphyxmini0504-figurelabel}
  \lean{PhyXMiniProblems.ProblemPhyXMini0504.FigureLabel}
  Mathematical labels printed in the primary image
\end{definition}
\begin{proof}
  This is the declared model definition of Figure Label.
\end{proof}
\begin{definition}[Figure Duration Mark]
  \label{def:physics:phyx-mini-0504:phyxminiproblems-problemphyxmini0504-figuredurationmark}
  \lean{PhyXMiniProblems.ProblemPhyXMini0504.FigureDurationMark}
  The two time intervals marked by double-headed arrows in the figure
\end{definition}
\begin{proof}
  This is the declared model definition of Figure Duration Mark.
\end{proof}
\begin{definition}[Pulse Train Figure]
  \label{def:physics:phyx-mini-0504:phyxminiproblems-problemphyxmini0504-pulsetrainfigure}
  \lean{PhyXMiniProblems.ProblemPhyXMini0504.PulseTrainFigure}
  \uses{def:physics:phyx-mini-0504:phyxminiproblems-problemphyxmini0504-durationquantity, def:physics:phyx-mini-0504:phyxminiproblems-problemphyxmini0504-pulsetrainshape, def:physics:phyx-mini-0504:phyxminiproblems-problemphyxmini0504-figurelabelrole, def:physics:phyx-mini-0504:phyxminiproblems-problemphyxmini0504-figurelabel, def:physics:phyx-mini-0504:phyxminiproblems-problemphyxmini0504-figuredurationmark}
  The labelled rectangular-pulse figure and its dimensionful arrow data
\end{definition}
\begin{proof}
  This is the declared model definition of Pulse Train Figure.
\end{proof}
\begin{definition}[Bandwidth Limited Pulse Train]
  \label{def:physics:phyx-mini-0504:phyxminiproblems-problemphyxmini0504-bandwidthlimitedpulsetrain}
  \lean{PhyXMiniProblems.ProblemPhyXMini0504.BandwidthLimitedPulseTrain}
  \uses{def:physics:phyx-mini-0504:phyxminiproblems-problemphyxmini0504-durationquantity, def:physics:phyx-mini-0504:phyxminiproblems-problemphyxmini0504-frequencyquantity, def:physics:phyx-mini-0504:phyxminiproblems-problemphyxmini0504-pulsetrainfigure}
  The independent physical quantities in the bandwidth-limited pulse system. The maximum rate is an independent inverse-time quantity. Its relation to the independent period, pulse duration, and bandwidth is imposed only by the general laws below; no numerical answer is built into this structure
\end{definition}
\begin{proof}
  This is the declared model definition of Bandwidth Limited Pulse Train.
\end{proof}
\begin{definition}[Matches Problem Statement And Figure]
  \label{def:physics:phyx-mini-0504:phyxminiproblems-problemphyxmini0504-matchesproblemstatementandfigure}
  \lean{PhyXMiniProblems.ProblemPhyXMini0504.MatchesProblemStatementAndFigure}
  \uses{def:physics:phyx-mini-0504:phyxminiproblems-problemphyxmini0504-durationreadout, def:physics:phyx-mini-0504:phyxminiproblems-problemphyxmini0504-bandwidthinkilohertz, def:physics:phyx-mini-0504:phyxminiproblems-problemphyxmini0504-bandwidthlimitedpulsetrain}
  Data transcribed from the prose and the primary image: a periodic rectangular pulse train, horizontal time axis t, pulse-width arrow Delta t, period arrow T = 2 Delta t, and 200 kHz electronics bandwidth. The last-but-one field is the semantic content of the period label, stated in every time unit. No maximum-rate value or answer choice occurs here
\end{definition}
\begin{proof}
  This is the declared model definition of Matches Problem Statement And Figure.
\end{proof}
\begin{definition}[Has Physical Pulse Train Parameters]
  \label{def:physics:phyx-mini-0504:phyxminiproblems-problemphyxmini0504-hasphysicalpulsetrainparameters}
  \lean{PhyXMiniProblems.ProblemPhyXMini0504.HasPhysicalPulseTrainParameters}
  \uses{def:physics:phyx-mini-0504:phyxminiproblems-problemphyxmini0504-durationinseconds, def:physics:phyx-mini-0504:phyxminiproblems-problemphyxmini0504-frequencyinhertz, def:physics:phyx-mini-0504:phyxminiproblems-problemphyxmini0504-pulserateinpulsespersecond, def:physics:phyx-mini-0504:phyxminiproblems-problemphyxmini0504-bandwidthlimitedpulsetrain}
  Positivity conditions for a nondegenerate pulse train and electronics system
\end{definition}
\begin{proof}
  This is the declared model definition of Has Physical Pulse Train Parameters.
\end{proof}
\begin{definition}[Satisfies Ideal Pulse Bandwidth Laws]
  \label{def:physics:phyx-mini-0504:phyxminiproblems-problemphyxmini0504-satisfiesidealpulsebandwidthlaws}
  \lean{PhyXMiniProblems.ProblemPhyXMini0504.SatisfiesIdealPulseBandwidthLaws}
  \uses{def:physics:phyx-mini-0504:phyxminiproblems-problemphyxmini0504-durationreadout, def:physics:phyx-mini-0504:phyxminiproblems-problemphyxmini0504-frequencyreadout, def:physics:phyx-mini-0504:phyxminiproblems-problemphyxmini0504-bandwidthlimitedpulsetrain}
  The ideal bandwidth-limited resolution model used by the numerical problem. At the limiting pulse width the reciprocal bandwidth-duration relation is B Delta-t = 1. A periodic train sends one pulse per start-to-start period, so its limiting pulse rate obeys R T = 1. Both laws are expressed in every time unit and contain neither 100000 pulses/s nor answer label C
\end{definition}
\begin{proof}
  This is the declared model definition of Satisfies Ideal Pulse Bandwidth Laws.
\end{proof}
\begin{definition}[Answer Choice]
  \label{def:physics:phyx-mini-0504:phyxminiproblems-problemphyxmini0504-answerchoice}
  \lean{PhyXMiniProblems.ProblemPhyXMini0504.AnswerChoice}
  Labels of the four multiple-choice answers
\end{definition}
\begin{proof}
  This is the declared model definition of Answer Choice.
\end{proof}
\begin{definition}[displayed Pulse Rate In Pulses Per Second]
  \label{def:physics:phyx-mini-0504:phyxminiproblems-problemphyxmini0504-displayedpulserateinpulsespersecond}
  \lean{PhyXMiniProblems.ProblemPhyXMini0504.displayedPulseRateInPulsesPerSecond}
  \uses{def:physics:phyx-mini-0504:phyxminiproblems-problemphyxmini0504-answerchoice}
  Pulse-rate number printed beside each choice, in pulses per second
\end{definition}
\begin{proof}
  This is the declared model definition of displayed Pulse Rate In Pulses Per Second.
\end{proof}
\begin{definition}[recorded Dataset Answer]
  \label{def:physics:phyx-mini-0504:phyxminiproblems-problemphyxmini0504-recordeddatasetanswer}
  \lean{PhyXMiniProblems.ProblemPhyXMini0504.recordedDatasetAnswer}
  \uses{def:physics:phyx-mini-0504:phyxminiproblems-problemphyxmini0504-answerchoice}
  The answer label recorded in the source dataset
\end{definition}
\begin{proof}
  This is the declared model definition of recorded Dataset Answer.
\end{proof}
\begin{definition}[Matches Answer Choice]
  \label{def:physics:phyx-mini-0504:phyxminiproblems-problemphyxmini0504-matchesanswerchoice}
  \lean{PhyXMiniProblems.ProblemPhyXMini0504.MatchesAnswerChoice}
  \uses{def:physics:phyx-mini-0504:phyxminiproblems-problemphyxmini0504-pulserateinpulsespersecond, def:physics:phyx-mini-0504:phyxminiproblems-problemphyxmini0504-bandwidthlimitedpulsetrain, def:physics:phyx-mini-0504:phyxminiproblems-problemphyxmini0504-answerchoice, def:physics:phyx-mini-0504:phyxminiproblems-problemphyxmini0504-displayedpulserateinpulsespersecond}
  A physical maximum pulse rate agrees with the value printed for a choice
\end{definition}
\begin{proof}
  This is the declared model definition of Matches Answer Choice.
\end{proof}
\begin{lemma}[maximum pulse rate is half bandwidth]
  \label{lem:physics:phyx-mini-0504:phyxminiproblems-problemphyxmini0504-maximum-pulse-rate-is-half-bandwidth}
  \lean{PhyXMiniProblems.ProblemPhyXMini0504.maximum_pulse_rate_is_half_bandwidth}
  \uses{def:physics:phyx-mini-0504:phyxminiproblems-problemphyxmini0504-frequencyreadout, def:physics:phyx-mini-0504:phyxminiproblems-problemphyxmini0504-bandwidthlimitedpulsetrain, def:physics:phyx-mini-0504:phyxminiproblems-problemphyxmini0504-matchesproblemstatementandfigure, def:physics:phyx-mini-0504:phyxminiproblems-problemphyxmini0504-hasphysicalpulsetrainparameters, def:physics:phyx-mini-0504:phyxminiproblems-problemphyxmini0504-satisfiesidealpulsebandwidthlaws}
  For an ideal pulse train with T = 2 Delta-t, the limiting reciprocal laws give a maximum rate equal to one half of the system bandwidth. This is a derived relation, not a law field or setup definition
\end{lemma}
\begin{proof}
  The conclusion follows from the governing relations and figure readouts named in the statement of maximum pulse rate is half bandwidth.
\end{proof}
% --- Archon declaration topology end ---
% --- Archon physics formalization source end ---
